\section{Wohinüber}
\stepcounter{ctrpoeminyear}
%\addcontentsline{toc}{section}{\thechapter.\arabic{ctrpoeminyear}\; Wohinüber}

\begin{strophe}
\vers{Der Zug ist abgefahren.}
\vers{Hier steh ich nun}
\vers{ganz allein vor leeren Schienen,}
\vers{vor leeren Hülsen falscher Gäste,}
\vers{die weinen über einen Staub,}
\vers{die nicht ahnen was dort macht.}
\vers{Kein Zweifel, jetzt bin ich frei.}
\vers{Frei von Liebe, frei von Lust.}
\vers{Gefangen in Problemen und Todeswünschen,}
\vers{gefangen in meinem Kopf im Todesakt.}
\vers{Tot und wunderschön tot.}
\vers{Kein Traum, der mich träumen lässt.}
\vers{Kein Wein, der mich lachen lässt.}
\vers{Keine Hoffnung innerhalb der verlorenen Welt.}
\vers{Alles ist vorbei.}
\vers{Konntest du dich nicht benehmen}
\vers{wenn du solltest aufrecht stehen?}
\vers{So schlägt auf dich die Peitsche}
\vers{der Realität bis du bleich wirst.}
\vers{Bleich wie weiß wie Schnee wie Tod.}
\end{strophe}

03.02.2021, minimal überarbeitet am 06.02.2021

%Christoph Koloska
